\section{Changing the domain and composing the constructions}
\label{sec:transfer}

\subsection{A finite partition theorem on metric spaces}

The cumulative probes exploit order on the line. A simpler reconstruction
works on any bounded rational presentation with an explicitly supplied
partition of unity. This gives a comparison for the precision saved by the
one-dimensional construction and identifies the data needed on a new domain.

Let $D$ be a rational metric basis of diameter at most $H>0$, and let an
extensional test language contain constants, rational linear combinations and
the following supplied test functions. At level $n$ produce a nonempty list of
sites $x_1,\ldots,x_m$ and nonnegative tests $\psi_i$ with
$\sum_i\psi_i=1$ and $\psi_i(x)>0\Rightarrow d(x,x_i)\le h_n$.
Both the producer and the support and partition proofs are input contracts.
Assume also that the language contains the functions
$\min_i(v_i+d(x,x_i))$ for finite rational lists of sites and ordinates.
Distance evaluations and the displayed test constructors must be PR on basis
codes. Order claims about tests are proved pointwise on the dense basis.
In particular the positivity rule for $L$ applies to such proved comparisons.
Use a supplied schedule making $h_n$ small. If $L$ is presented on a smaller
test language, the additional constructors and their properties are obligations.

\begin{theorem}[Reconstruction from a partition]
Given a positive normalized rational-linear evaluator $L$, read each
$L(\psi_i)$ with error $\delta=2^{-r}$, where $m\delta\le1/2$.
In $O(m)$ rational operations after the oracle calls one can produce a
finite probability law $P$ on the sites satisfying
\[
             |P(f)-L(f)|\le K(h_n+Hm\delta)
\]
for every admitted $K$-Lipschitz test $f$. There is a PR procedure to compute
the maximum discrepancy of two finite laws against all 1-Lipschitz tests
in this language.
\end{theorem}
\begin{proof}
Write $c_i=L(\psi_i)$. Truncate negative readings to zero to obtain rational
$a_i$ with $|a_i-c_i|\le\delta$, and set $s=\sum_i a_i$.
Normalization and linearity give $|s-1|\le m\delta$. In particular
$s\ge1/2$. Set $P_i=a_i/s$. The sum of the normalization corrections
$\sum_i|P_i-a_i|$ is $|1-s|$. Thus $\sum_i|P_i-c_i|\le2m\delta$.
Here absolute values on named scalars can equally be evaluated by the regular
absolute-value operation. For an argument using only rational sign tests,
compare to arbitrary rational approximants to the $c_i$ and remove the
accumulated slack as in Section~\ref{sec:laws}.

The positive interpolant $Tf=\sum_i f(x_i)\psi_i$ approximates $f$ within
$Kh_n$ on the basis. Indeed its terms have nonnegative weights summing to
one and the support condition bounds every nonzero discrepancy. This can also
be proved without deciding equality of a named weight: check strict failure
of the rational distance test; in that case nonnegative $\psi_i$ with no
positive evidence is zero by the rational gap rule. Subtract the midpoint of
the greatest and least of the rational nodal ordinates when $f$ is a rational
test on the given sites. The deviations at the sites are at most $KH/2$,
and subtracting a constant leaves the mass discrepancy unchanged. This
proves the bound on replacing $c_i$ by $P_i$. For named nodal values choose rational $u_i$ within $e$ and use the admitted test $T_u=\sum_i u_i\psi_i$. Its uniform discrepancy from $f$ is at most $Kh_n+e$. Centering the rational ordinates bounds their magnitude by $KH/2+e$, and replacing the target ordinates in $P(f)$ adds at most $e$. Remove the resulting slack by the rational gap rule. Only the rational masses enter the producer.

Here is an explicit procedure for the second assertion. On the union of the
two supports fix a base vertex and normalize a potential to zero there.
Let $B$ be a common positive denominator for all pairwise distances. Enumerate
integer ordinates $v_i$ bounded in magnitude by $\lceil BH\rceil$, with
$v_0=0$ and $|v_i-v_j|\le B d(x_i,x_j)$. Return the largest signed score
of the ordinates $v_i/B$ against the difference of the masses. Each array
extends with these same nodal values by the finite minimum stipulated above.
This test is 1-Lipschitz by the triangle inequality.

For completeness of the search, start with any normalized \emph{rational}
feasible array $y$. For $0\le t<1$ set $v_i(t)=\lfloor B y_i+t\rfloor$.
Because $Bd(x_i,x_j)$ is an integer, these rounded arrays are still feasible,
and $v_0(t)=0$. The average of $v_i(t)$ over $t$ is $B y_i$, as one sees
by splitting at the single jump of that coordinate. Refinement by all jump
locations expresses the original score as a finite rational convex combination
of enumerated scores. The largest score dominates this combination. For a
test with named nodal values, the inequality follows by approximation: choose
rational approximations $u_i$ within $e$ and use the rational potential
$y_j=\min_i(u_i+d(x_i,x_j))$. It is feasible and within $e$ of the original
at each site. Normalize at the base vertex; the constant cancels, and the
score error is at most $2e$. Remove this slack. Reversing the test's sign
also bounds the absolute discrepancy.
\end{proof}

After merging zero-distance sites it is a metric: to detect a mass discrepancy at a single site, give that site the positive value of its least distance to the other sites and give all other sites value zero. The extension above is feasible. A support with one site has only one unit law. This test distance yields a regular law name by the same maximizing-test
two-run argument as on the line. Refine until $h_n\le2^{-p-1}$ and
$Hm\delta\le2^{-p-1}$, keeping $m\delta\le1/2$. The finite dual
enumeration is a decidability proof and a general algorithm; its worst-case
search size is exponential in the support size even for a fixed ordinate
range greater than one. The ordered scan in Lemma~\ref{lem:dual} removes
this search on the line and supplies a sparse primal coupling there.
The theorem above establishes a metric of continuous tests directly.

For example, on $[0,1]^d$ in the maximum distance use tensor products of
ordinary linear hats on the dyadic grid. Their products are nonnegative,
sum to one by multiplying the $d$ one-dimensional identities, and are
supported within one mesh step. With fixed precision $n$, $m=(2^n+1)^d$,
$H=1$ and $h_n=2^{-n}$. Rational piecewise polynomial tests with min and
max constructors suffice for these expressions; witnesses of positivity
can be hypotheses or explicit algebraic proofs. This instance needs no
general positivity decider for all expressions. For this simple normalization
method the required raw query precision grows as $\lceil\log_2 m\rceil$.
On $[0,1]$, Section~\ref{sec:laws} eliminates that term by probing the
cumulative sums directly.

\subsection{A quantile from a transport name}

The operator for roots accepts a general function evaluator. The law
construction can provide it, even though a distribution function at an
atom was a counterexample. Suppose now the inputs have the following density bounds on rational test functions
\[
 2^{-b}\int_0^1 h\ \le L(h)\le 2^d\int_0^1 h
 \qquad(0\le h\in\PL)
\]
with supplied integers $b,d$. Every displayed polygonal integral is a finite
rational trapezoid sum. For a rational $x\in[0,1]$ let $R_{x,\eta}$ be the
right ramp of width $\eta$, equal to one through $x$ and zero past $x+\eta$,
restricted to $[0,1]$. Its area lies between $x$ and $x+\eta$. At fixed
$x$ the ramps are ordered by width. The upper density bound shows that the
two expectations at widths $\eta,\eta'$ differ by at most
$2^d\max(\eta,\eta')$. Thus they define a scalar $C(x)$ and the error at
one width is at most $2^d\eta$. For $y\ge x$ the difference of the two
equal-width ramps is nonnegative and its integral tends to $y-x$ with
error at most $\eta$. Applying both density comparisons to that difference
and removing the auxiliary width proves
\[
 2^{-b}(y-x)\le C(y)-C(x)\le2^d(y-x).
\]
The resulting function has endpoint values $C(0)=0,C(1)=1$. With
a semantic probability measure, $C$ is its distribution function.

Given a law name $M$, to evaluate $C(x)$ at precision $p$ use
$\eta=2^{-(p+d+2)}$ and the single finite law at level $\ell=2p+d+4$.
The ramp has Lipschitz constant at most $1/\eta$. Its value under $M_\ell$
has discrepancy at most $2^{-\ell}/\eta=e_p/4$ from $L(R_{x,\eta})$.
The width error is also at most $e_p/4$. Both estimates include $x=0,1$.
Subtract a rational target mass $\alpha\in[0,1]$ to obtain observations
for Section~\ref{sec:conditioned-zero} with $g(k)=k+b$ and $\mu(k)=k+d$.
Consequently one finite-law observation per halving produces a coherent name
for the unique $\alpha$-quantile. At halving $j$ its level is
$2(j+b+3)+d+4$. This counts law observations. The number of atoms encoded
by one observed law still affects the work to integrate the ramp. The upper
density promise supplies a CDF evaluator; the lower promise supplies the
inverse modulus. Their roles are distinct. The construction feeds its
explicit moduli through the general composition rule instead of taking an
inverse of a black-box distribution function.


For an exact instance, take density $1/2+x$ on $[0,1]$. Here $b=d=1$
suffice, and the distribution function for independent checking is
$(x+x^2)/2$. The program \texttt{examples/quantiles.py} represents the
level-$\ell$ law by moving each dyadic cell's mass to its left endpoint.
This costs at most $2^{-\ell}$. It sums a ramp on this law from the first
two prefix moments, so it can use very fine levels without materializing the
grid. At $\alpha=7/15$, 32 calls give the bracket
$[5049260259,5049260261]/8589934592$. Substitution in the quadratic verifies
the inequalities at both endpoints. The program also exercises $\alpha=0$
and $\alpha=1$, and checks the compressed evaluation against direct finite sums.
